\AtBeginSection{
% \begin{frame}
% \frametitle{\kap. Beispiel Kugelvolumen}
% \tableofcontents[current]
% \end{frame}
}
%%%%%%%%%%%%%%%%%%%%%%%%%%%%%%%%%%%%%%%%%%%%%%%%%%%%%%%%%%%%%%%%%%%%%%%%
\section{Ausblick}
\begin{frame}
\frametitle{\kap. Ausblick}%
\tableofcontents
\end{frame}
%%%%%%%%%%%%%%%%%%%%%%%%%%%%%%%%%%%%%%%%%%%%%%%%%%%%%%%%%%%%%%%%%%%%%%%%
\def\stitle{Guter Start ins Praktikum}%
\subsection{\stitle}\label{S:Guter Start}
\begin{frame}[t]%
  \frametitle{\kap.\ref{S:Guter Start} \stitle}%
\medskip

In dieser \"Ubung wird anhand eines Beispiels beschrieben wie ein Java-Programm \"ubersichtlich gestaltet werden kann.

\pause
\begin{itemize}
    \item Wichtig f\"ur eigenes Verst\"andnis des Programms
    \item Wichtig f\"ur Dokumentation des Programms
    \item Wichtig f\"ur das Erhalten von Testaten
    \item Unwichtig f\"ur den Computer
\end{itemize}

\pause \bigskip
\begin{block}{Programmierstil}
Einr\"ucken, Kommentare, Klammerung, Namen
\end{block}

\pause \bigskip
\begin{block}{Hinweis}
Aufgabe 1 (Pflichtaufgabe) und 2 lassen sich auf \"ahnliche Weise l\"osen
\end{block}

\end{frame}

\def\stitle{Beispiel Obelisk}%
\subsection{\stitle}\label{S:Beispiel Obelisk}
\begin{frame}[t]%
    \frametitle{\kap.\ref{S:Beispiel Obelisk} \stitle}%
\medskip

\begin{itemize}
 \item Ein \emph{Obelisk} ist ein Polyeder mit trapezf\"ormigen Seitenfl\"achen
 \item Es seien
 \begin{itemize}
  \item $a$ und $b$ die Seitenl\"angen der Grundfl\"ache
  \item $c$ und $d$ die Seitenl\"angen der Deckfl\"ache
  \item $h$ die H\"ohe des Obelisken
 \end{itemize}
\end{itemize}

\begin{center}
\includegraphics[width=0.5\textwidth]{guterStartInsPraktikum/bilder/obelisk}
\end{center}

%  \item Schreiben Sie ein Java-Programm, welches das Kugelvolumen einer beliebigen Kugel berechnet
% \end{itemize}
% \smallskip
%
% Ben"otigte Werkzeuge:
% \begin{itemize}
%  \item Standard Programmierger"ust
%  \item Geeignete Variablen f"ur $r$ und $v$
%  \item Einlesen von Variablen
%  \item Wert von $\pi$
%  \item Ausgeben des Ergebnisses
% \end{itemize}
\end{frame}

\begin{frame}[t]%
    \frametitle{\kap.\ref{S:Beispiel Obelisk} \stitle}%
\medskip
Dann ist:
\begin{itemize}
\item $G_1 = ab$ der Fl\"acheninhalt der Grundfl\"ache.
\item $G_2 = cd$ der Fl\"acheninhalt der Deckfl\"ache.
\item $\displaystyle A_1 = \frac{a+c}{2} \sqrt{h^2 + \Bigl(\frac{b-d}{2}\Bigr)^2}$
der Fl\"acheninhalt zweier Mantelfl\"achen.
\item $\displaystyle A_2 = \frac{b+d}{2} \sqrt{h^2 + \Bigl(\frac{a-c}{2}\Bigr)^2}$
der Fl\"acheninhalt der anderen beiden Mantelfl\"achen.
\item $M = 2 A_1 + 2 A_2$ der Fl\"acheninhalt der gesamten Mantelfl\"ache.
\item $O = M + G_1 + G_2$ der Fl\"acheninhalt der Oberfl\"ache .
\item $\displaystyle V = \frac{h}{6}\bigl(G_1 + (a+c)(b+d) + G_2\bigr)$ das Volumen.
\end{itemize}
\medskip

Schreiben Sie ein Java--Programm, welches die Seitenl\"angen $a$, $b$, $c$ und $d$ sowie
die H\"ohe $h$ eines Obelisken von der Konsole einliest und anschlie\ss end den Fl\"acheninhalt
der Mantelfl\"ache $M$, den der Oberfl\"ache $O$ sowie das Volumen $V$ ausgibt.

\end{frame}

\begin{frame}[t]%
    \frametitle{\kap.\ref{S:Beispiel Obelisk} \stitle}%
\medskip


\codeblock{\noindent
\cmt{
/* Berechnung von\\
\ * \ Mantelflaeche M\\
\ * \ Oberflaeche O\\
\ * \ Volumen V\\
\ * eines Obelisken\\
\ * \\
\ * Eingabe: a,b,c,d (Seitenlaengen) h (Hoehe)\\
\ * Ausgabe: M, O, V\\
\ * Author: Stefan Findeisen 2016\\
\ */\\[1em]
 }
public class Obelisk \{\\
\tab public static void main(String[] args) \{\\
\tab \tab \cmt{/* Hier steht das Hauptprogramm */}\\
\tab \tab ...\\
\tab \}\\
\}\\
}
\end{frame}

\begin{frame}[t]%
    \frametitle{\kap.\ref{S:Beispiel Obelisk} \stitle}%
\medskip

\codeblock{\noindent
\cmt{
/* Berechnung von\\
\ * \ Mantelflaeche M\\
\ * \ Oberflaeche O\\
\ * \ Volumen V\\
\ * eines Obelisken\\
\ * \\
\ * Eingabe: a,b,c,d (Seitenlaengen) h (Hoehe)\\
\ * Ausgabe: M, O, V\\
\ * Author: Stefan Findeisen 2016\\
\ */\\[1em]
 }
public class Obelisk \{\\
\tab public static void main(String[] args) \{\\
\tab \tab \cmt{/* Hier steht das Hauptprogramm */}\\
\tab \tab \cmt{/* Setze die Standardeinstellungen */}\\
\tab \tab {\color{red} Locale.setDefault(Locale.US);}\\
\tab \tab ...\\
\tab \}\\
\}\\
}
\end{frame}

\begin{frame}[t]%
    \frametitle{\kap.\ref{S:Beispiel Obelisk} \stitle}%
\medskip

\codeblock{\noindent
\cmt{
\ ...\\
\ * Eingabe: a,b,c,d (Seitenlaengen) h (Hoehe)\\
\ * Ausgabe: M, O, V\\
\ * Author: Stefan Findeisen 2016\\
\ */\\[1em]
 }
public class Obelisk \{\\
\tab public static void main(String[] args) \{\\
\tab \tab \cmt{/* Hier steht das Hauptprogramm */}\\
\tab \tab \cmt{/* Setze die Standardeinstellungen */}\\
\tab \tab Locale.setDefault(Locale.US);\\[1em]
\tab \tab \cmt{/* Daten des Obelisken */}\\
\tab \tab {\color{red} double a, b, c, d;} \cmt{// Seitenlaengen}\\
\tab \tab {\color{red} double h;} \cmt{ // Hoehe}\\
\tab \tab ...\\
\tab \}\\
\}\\
}
\end{frame}

\begin{frame}[t]%
    \frametitle{\kap.\ref{S:Beispiel Obelisk} \stitle}%
\medskip

\codeblock{\noindent
\cmt{...}\\[1em]
public class Obelisk \{\\
\tab public static void main(String[] args) \{\\
\tab \tab \cmt{/* Hier steht das Hauptprogramm */}\\
\tab \tab \cmt{/* Setze die Standardeinstellungen */}\\
\tab \tab Locale.setDefault(Locale.US);\\[1em]
\tab \tab \cmt{/* Daten des Obelisken */}\\
\tab \tab {\color{red} double a, b, c, d;} \cmt{// Seitenlaengen}\\
\tab \tab {\color{red} double h;} \cmt{ // Hoehe}\\[1em]
\tab \tab \cmt{/* Eingabe */}\\
\tab \tab \cmt{/* Hilfstext: */}\\
\tab \tab {\color{red} System.out.println(''Geben Sie die Seitenlaengen (a,b,c,d) '' +\\
\tab \tab \tab         ''sowie die Hoehe h eines Obelisken ein:'');}\\
\tab \tab ...\\
\tab \}\\
\}\\
}
\end{frame}

\begin{frame}[t]%
    \frametitle{\kap.\ref{S:Beispiel Obelisk} \stitle}%
\vspace{-1em}
\codeblock{\noindent
public class Obelisk \{\\
\tab public static void main(String[] args) \{\\
\tab \tab ...\\
\tab \tab \cmt{/* Eingabe */}\\
\tab \tab \cmt{/* Hilfstext: */}\\
\tab \tab System.out.println(''Geben Sie die Seitenlaengen (a,b,c,d) '' +\\
\tab \tab \tab         ''sowie die Hoehe h eines Obelisken ein:'');\\
\tab \tab \cmt{/* Eigentliche Eingabe mit Hilfstexten */}\\
\tab \tab {\color{red} Scanner sc = new Scanner(System.in);\\
\tab \tab   System.out.print(''a = '');\\
\tab \tab   a = sc.nextDouble();\\
\tab \tab   System.out.print(''b = '');\\
\tab \tab   b = sc.nextDouble();\\
\tab \tab   System.out.print(''c = '');\\
\tab \tab   c = sc.nextDouble();\\
\tab \tab   System.out.print(''d = '');\\
\tab \tab   d = sc.nextDouble();\\
\tab \tab   System.out.print(''h = '');\\
\tab \tab   h = sc.nextDouble();}\\
\tab \tab ...\\
\tab \}\\
\}\\
}
\end{frame}

\begin{frame}[t]%
    \frametitle{\kap.\ref{S:Beispiel Obelisk} \stitle}%

\codeblock{\noindent
{\color{red} import java.util.*;}  \cmt{// wird wegen ''Scanner'' benoetigt}\\[1em]
public class Obelisk \{\\
\tab public static void main(String[] args) \{\\
\tab \tab ...\\
\tab \tab \cmt{/* Eingabe */}\\
\tab \tab \cmt{/* Hilfstext: */}\\
\tab \tab System.out.println(''Geben Sie die Seitenlaengen (a,b,c,d) '' +\\
\tab \tab \tab         ''sowie die Hoehe h eines Obelisken ein:'');\\
\tab \tab \cmt{/* Eigentliche Eingabe mit Hilfstexten */}\\
\tab \tab {\color{red} Scanner sc = new Scanner(System.in);\\
\tab \tab   ...}\\
\tab \tab ...\\
\tab \}\\
\}\\
}
\end{frame}

\begin{frame}[t]%
    \frametitle{\kap.\ref{S:Beispiel Obelisk} \stitle}%

\codeblock{\noindent
import java.util.*;  \cmt{// wird wegen ''Scanner'' benoetigt}\\[1em]
public class Obelisk \{\\
\tab public static void main(String[] args) \{\\
\tab \tab ...\\
\tab \tab \cmt{/* Berechung */}\\
}
\begin{itemize}
\item $G_1 = ab$ Fl\"acheninhalt der Grundfl\"ache:
\end{itemize}
\codeblock{\noindent
\tab \tab {\color{red} double G1 = a * b;}\\
}
\begin{itemize}
\item $G_2 = cd$ Fl\"acheninhalt der Deckfl\"ache:
\end{itemize}
\codeblock{\noindent
\tab \tab {\color{red} double G2 = c * d;}\\
\tab \tab ...\\
\tab \}\\
\}\\
}
\end{frame}

\begin{frame}[t]%
    \frametitle{\kap.\ref{S:Beispiel Obelisk} \stitle}%

\codeblock{\noindent
import java.util.*;  \cmt{// wird wegen ''Scanner'' benoetigt}\\[1em]
public class Obelisk \{\\
\tab public static void main(String[] args) \{\\
\tab \tab ...\\
\tab \tab \cmt{/* Berechung */}\\
\tab \tab ...\\
}
\begin{itemize}
\item $\displaystyle A_1 = \frac{a+c}{2} \sqrt{h^2 + \Bigl(\frac{b-d}{2}\Bigr)^2}$
Fl\"acheninhalt zweier Mantelfl\"achen:
\end{itemize}
\codeblock{\noindent
\tab \tab {\color{red} double s1 = 0.5 * (a + c);\\
\tab \tab double t1 = 0.5 * (b - d);}\\
\tab \tab double A1 = s1 * Math.sqrt(h * h + t1 * t1);\\
}
\begin{itemize}
\item $\displaystyle A_2 = \frac{b+d}{2} \sqrt{h^2 + \Bigl(\frac{a-c}{2}\Bigr)^2}$
Fl\"acheninhalt der anderen Mantelfl\"achen:
\end{itemize}
\codeblock{\noindent
\tab \tab {\color{red} double s2 = 0.5 * (b + d);\\
\tab \tab double t2 = 0.5 * (a - c);}\\
\tab \tab double A2 = s2 * Math.sqrt(h * h + t2 * t2);\\
\tab \tab ...\\
\tab \}\\
\}\\
}
\end{frame}

\begin{frame}[t]%
    \frametitle{\kap.\ref{S:Beispiel Obelisk} \stitle}%

\codeblock{\noindent
import java.util.*;  \cmt{// wird wegen ''Scanner'' benoetigt}\\[1em]
public class Obelisk \{\\
\tab public static void main(String[] args) \{\\
\tab \tab ...\\
\tab \tab \cmt{/* Berechung */}\\
\tab \tab ...\\
}
\begin{itemize}
\item $\displaystyle A_1 = \frac{a+c}{2} \sqrt{h^2 + \Bigl(\frac{b-d}{2}\Bigr)^2}$
Fl\"acheninhalt zweier Mantelfl\"achen:
\end{itemize}
\codeblock{\noindent
\tab \tab double s1 = 0.5 * (a + c);\\
\tab \tab double t1 = 0.5 * (b - d);\\
\tab \tab double A1 = s1 * {\color{red} Math.sqrt}(h * h + t1 * t1);\\
}
\begin{itemize}
\item $\displaystyle A_2 = \frac{b+d}{2} \sqrt{h^2 + \Bigl(\frac{a-c}{2}\Bigr)^2}$
Fl\"acheninhalt der anderen Mantelfl\"achen:
\end{itemize}
\codeblock{\noindent
\tab \tab double s2 = 0.5 * (b + d);\\
\tab \tab double t2 = 0.5 * (a - c);\\
\tab \tab double A2 = s2 * {\color{red} Math.sqrt}(h * h + t2 * t2);\\
\tab \tab ...\\
\tab \}\\
\}\\
}
\end{frame}

\begin{frame}[t]%
    \frametitle{\kap.\ref{S:Beispiel Obelisk} \stitle}%

\codeblock{\noindent
import java.util.*;  \cmt{// wird wegen ''Scanner'' benoetigt}\\[1em]
public class Obelisk \{\\
\tab public static void main(String[] args) \{\\
\tab \tab ...\\
\tab \tab \cmt{/* Berechung */}\\
\tab \tab ...\\
}
\begin{itemize}
 \item $M = 2 A_1 + 2 A_2$ Fl\"acheninhalt der gesamten Mantelfl\"ache:
\end{itemize}
\codeblock{\noindent
\tab \tab {\color{red} double M = 2 * A1 + 2 * A2;}\\
}
\begin{itemize}
 \item $O = M + G_1 + G_2$ Fl\"acheninhalt der Oberfl\"ache:
\end{itemize}
\codeblock{\noindent
\tab \tab {\color{red} double O = M + G1 + G2;}\\
}
\begin{itemize}
 \item $\displaystyle V = \frac{h}{6}\bigl(G_1 + (a+c)(b+d) + G_2\bigr)$ Volumen:
\end{itemize}
\codeblock{\noindent
\tab \tab {\color{red} double V = h / 6 * (G1 + (a + c) * (b + d) + G2);}\\
\tab \tab ...\\
\tab \}\\
\}\\
}
\end{frame}

\begin{frame}[t]%
    \frametitle{\kap.\ref{S:Beispiel Obelisk} \stitle}%

\codeblock{\noindent
import java.util.*;  \cmt{// wird wegen ''Scanner'' benoetigt}\\[1em]
public class Obelisk \{\\
\tab public static void main(String[] args) \{\\
\tab \tab ...\\
\tab \tab \cmt{/* Ausgabe */}\\
\tab \tab {\color{red} System.out.println(''Mantelflaeche = '' + M);\\
\tab \tab System.out.println(''Oberflaeche = '' + O);\\
\tab \tab System.out.println(''Volumen = '' + V);}\\
\tab \}\\
\}\\
}
\end{frame}

\begin{frame}[t]%
    \frametitle{\kap.\ref{S:Beispiel Obelisk} \stitle}%
\tiny
\codeblock{\noindent
import java.util.*; \cmt{// wird wegen ''Scanner'' benoetigt}\\[1em]
\cmt{
/* Berechnung von Mantelflaeche M, Oberflaeche O, Volumen V, eines Obelisken\\
 * Eingabe: a,b,c,d (Seitenlaengen) h (Hoehe)\\
 * Ausgabe: M, O, V\\
 * Author: Stefan Findeisen 2016\\
 */}\\[1em]
public class Obelisk \{\\
\tab public static void main(String[] args) \{\\
\tab \tab \cmt{/* Hier steht das Hauptprogramm */}\\
\tab \tab \cmt{/* Setze die Standardeinstellungen */}\\
\tab \tab Locale.setDefault(Locale.US);\\[1em]
\tab \tab \cmt{/* Daten des Obelisken */}\\
\tab \tab double a, b, c, d; \cmt{// Seitenlaengen}\\
\tab \tab double h; \cmt{// Hoehe}\\
\tab \tab \cmt{/* Eingabe */\\
\tab \tab /* Hilfstext: */}\\
\tab \tab System.out.println(''Geben Sie die Seitenlaengen (a,b,c,d) '' +\\
\tab \tab \tab ''sowie die Hoehe h eines Obelisken ein:'');\\[1em]
\tab \tab \cmt{/* Eigentliche Eingabe mit Hilfstexten */}\\
\tab \tab Scanner sc = new Scanner(System.in);\\[1em]
\tab \tab System.out.print(''a = '');\\
\tab \tab a = sc.nextDouble();\\
\tab \tab System.out.print(''b = '');\\
\tab \tab b = sc.nextDouble();\\
\tab \tab System.out.print(''c = '');\\
\tab \tab c = sc.nextDouble();\\
\tab \tab System.out.print(''d = '');\\
\tab \tab d = sc.nextDouble();\\
\tab \tab System.out.print(''h = '');\\
\tab \tab h = sc.nextDouble();\\
\tab \tab ...\\
\tab \}\\
\}
}
\end{frame}

\begin{frame}[t]%
    \frametitle{\kap.\ref{S:Beispiel Obelisk} \stitle}%
\tiny
\codeblock{\noindent
import java.util.*; \cmt{// wird wegen ''Scanner'' benoetigt}\\[1em]
\cmt{
/* Berechnung von Mantelflaeche M, Oberflaeche O, Volumen V, eines Obelisken\\
 * Eingabe: a,b,c,d (Seitenlaengen) h (Hoehe)\\
 * Ausgabe: M, O, V\\
 * Author: Stefan Findeisen 2016\\
 */}\\[1em]
public class Obelisk \{\\
\tab public static void main(String[] args) \{\\
\tab \tab \cmt{/* Hier steht das Hauptprogramm */}\\
\tab \tab ...\\
\tab \tab \cmt{/* Berechung */}\\
\tab \tab double G1 = a * b;\\
\tab \tab double G2 = c * d;\\
\tab \tab double s1 = 0.5 * (a + c);\\
\tab \tab double s2 = 0.5 * (b + d);\\
\tab \tab double t1 = 0.5 * (b - d);\\
\tab \tab double t2 = 0.5 * (a - c);\\
\tab \tab double A1 = s1 * Math.sqrt(h * h + t1 * t1);\\
\tab \tab double A2 = s2 * Math.sqrt(h * h + t2 * t2);\\[1em]
\tab \tab double M = 2 * A1 + 2 * A2;\\
\tab \tab double O = M + G1 + G2;\\
\tab \tab double V = h / 6 * (G1 + (a + c) * (b + d) + G2);\\[1em]
\tab \tab \cmt{/* Ausgabe */}\\
\tab \tab System.out.println(''Mantelflaeche = '' + M);\\
\tab \tab System.out.println(''Oberflaeche = '' + O);\\
\tab \tab System.out.println(''Volumen = '' + V);\\
\tab \}\\
\}
}
\end{frame}
